\begin{figure}[htbp]
\centering
\begin{tikzpicture}
    \begin{axis}[
        width=10cm,
        height=10cm,
        xlabel={$x$},
        ylabel={$y$},
        xmin=-0.5, xmax=5,
        ymin=-0.5, ymax=5,
        axis equal image,
        xtick={0,1,2,3,4},
        ytick={0,1,2,3,4},
    ]
        % Constraint line: x + y = 4
        \addplot[
            very thick,
            UofSCAtlantic,
            domain=-0.5:5,
            samples=2,
        ] {4 - x};
        \node[above right, UofSCAtlantic, font=\small, rotate=-45] at (axis cs:1,3) {$g(x,y) = x + y - 4 = 0$};

        % Contours of objective f(x,y) = x^2 + y^2
        \addplot[
            contour lua={
                number=6,
                labels=false,
            },
            UofSCGarnet!50,
            domain=0:5,
            y domain=0:5,
            samples=40,
        ] {x^2 + y^2};

        % Contour labels
        \node[UofSCGarnet, font=\scriptsize] at (axis cs:1.5,0.3) {$f=4$};
        \node[UofSCGarnet, font=\scriptsize] at (axis cs:2.5,0.5) {$f=8$};
        \node[UofSCGarnet, font=\scriptsize] at (axis cs:3.2,1.0) {$f=12$};

        % Optimal point
        \addplot[only marks, mark=*, mark size=5pt, UofSCGarnet] coordinates {(2,2)};
        \node[above right, UofSCGarnet, font=\small\bfseries] at (axis cs:2.1,2.1) {$(2,2)$};

        % Gradient of f at optimal point (points outward from origin)
        \draw[->,UofSCGarnet,very thick] (axis cs:2,2) -- (axis cs:2.8,2.8);
        \node[above right, UofSCGarnet, font=\small] at (axis cs:2.8,2.8) {$\nabla f$};

        % Gradient of g at optimal point (perpendicular to constraint)
        \draw[->,UofSCAtlantic,very thick] (axis cs:2,2) -- (axis cs:2.5,2.5);
        \node[below right, UofSCAtlantic, font=\small] at (axis cs:2.5,2.5) {$\nabla g$};

        % Show they are parallel (key insight)
        \draw[<->,UofSCHorseshoe,dashed,thick] (axis cs:2.6,2.6) to[bend left=20] (axis cs:2.9,2.9);
        \node[right, UofSCHorseshoe, font=\scriptsize, align=left] at (axis cs:3.0,2.75) {$\nabla f \parallel \nabla g$\\(at optimum)};

        % Origin marker
        \addplot[only marks, mark=o, mark size=3pt, UofSC70Black] coordinates {(0,0)};
        \node[below right, font=\small] at (axis cs:0,0) {origin};

        % Unconstrained minimum (for comparison)
        \draw[dashed,UofSC50Black] (axis cs:0,0) -- (axis cs:2,2);
        \node[below, UofSC70Black, font=\scriptsize, align=center] at (axis cs:1,0.8) {distance\\$= 2\sqrt{2}$};

        % Feasible region label
        \node[UofSCAtlantic, font=\small, align=center] at (axis cs:0.8,4.2) {feasible set\\(constraint line)};

    \end{axis}
\end{tikzpicture}
\caption{Geometric interpretation of Lagrange multipliers. Minimizing $f(x,y) = x^2 + y^2$ (distance squared from origin) subject to $g(x,y) = x + y - 4 = 0$. Circular contours show level curves of $f$. At the constrained minimum $(2,2)$, the gradients $\nabla f$ and $\nabla g$ are parallel (both perpendicular to the constraint line). This is the geometric condition for optimality: we cannot improve the objective while staying on the constraint.}
\label{fig:constraint_visualization}
\end{figure}
